\documentclass[11pt]{article}

\usepackage[utf8]{inputenc}
\usepackage[margin=1in]{geometry}
\usepackage{amsmath, amssymb, amsthm}
\usepackage{CJKutf8}
\usepackage{enumitem}
\usepackage{hyperref}

\title{10-708 Midterm Review \\ Q9 \& Q10 讲解}
\author{}
\date{}

\begin{document}
\begin{CJK}{UTF8}{gbsn}
\maketitle

% ============================================================
\section{成绩分析}
% ============================================================

\subsection{成绩权重}

\begin{center}
\begin{tabular}{l|r}
  Component & Weight \\
  \hline
  Homework (4个) & 40\% \\
  Midterm & 25\% \\
  Final & 35\% \\
  Attendance (extra credit) & 3\% \\
\end{tabular}
\end{center}

\subsection{已知成绩}

\begin{center}
\begin{tabular}{l|r|r|r|l}
  & 你的分数 & Mean & Median & 备注 \\
  \hline
  HW1 & 103/105 = 98.1\% & 101.6 & 103.25 & 接近 median \\
  HW2 & 113.5/115 = 98.7\% & 109.25 & 111.5 & 高于 median \\
  HW3 & ? & & & 已交未出分 \\
  HW4 & ? & & & 未出 \\
  \hline
  Midterm & 29/70 = 41.4\% & 46.2 & 48.75 & 低于 mean 1.35 std \\
  \hline
  Attendance & $\sim$4/6 次 & & & $\approx 2\%$ extra credit \\
\end{tabular}
\end{center}

\subsection{精确计算}

假设 4 个 HW 等权，每个占 $40\% / 4 = 10\%$。HW1+HW2 已锁定，HW3+HW4 假设保持 $\sim$97\%。

已锁定 + 估计 + extra credit：
\begin{align}
  \text{HW1} &= \frac{103}{105} \times 10\% = 9.81\% \notag \\
  \text{HW2} &= \frac{113.5}{115} \times 10\% = 9.87\% \notag \\
  \text{HW3+HW4（估计）} &= 0.97 \times 10\% \times 2 = 19.4\% \notag \\
  \text{Midterm} &= \frac{29}{70} \times 25\% = 10.36\% \notag \\
  \text{Attendance} &\approx 2\% \notag \\
  \text{小计} &= \mathbf{53.44\%} \notag
\end{align}

\subsection{期末需要多少分？}

Prof 说只有 curve down 没有 curve up --- 不指望 curve。按 CMU 标准 cutoff：

\begin{center}
\begin{tabular}{l|c|c}
  目标 & Final 需要 & 备注 \\
  \hline
  \textbf{B} ($\geq 83\%$) & $\frac{83 - 53.44}{35} = \mathbf{84.5\%}$ & 难但有可能 \\[4pt]
  \textbf{B$-$} ($\geq 80\%$) & $\frac{80 - 53.44}{35} = \mathbf{75.9\%}$ & 可行 \\[4pt]
  \textbf{C$+$} ($\geq 77\%$) & $\frac{77 - 53.44}{35} = \mathbf{67.3\%}$ & 稳的 \\
\end{tabular}
\end{center}

\subsection{学期 GPA 影响}

另外两门课：A+ (4.33) 和 A (4.0)。三门课等学分。

\begin{center}
\begin{tabular}{l|c|c}
  10-708 成绩 & GPA 点数 & 学期 GPA \\
  \hline
  B & 3.33 & $(4.33 + 4.0 + 3.33)/3 = \mathbf{3.89}$ \\
  B$-$ & 3.0 & $(4.33 + 4.0 + 3.0)/3 = \mathbf{3.78}$ \\
  C$+$ & 2.67 & $(4.33 + 4.0 + 2.67)/3 = \mathbf{3.67}$ \\
  C & 2.33 & $(4.33 + 4.0 + 2.33)/3 = \mathbf{3.55}$ \\
\end{tabular}
\end{center}

\vspace{0.3em}
\noindent\fbox{\parbox{\dimexpr\textwidth-2\fboxsep-2\fboxrule}{%
\textbf{结论。}期末 76\% → B$-$（学期 GPA 3.78），85\% → B（学期 GPA 3.89）。HW 98\%+ 是安全网，另外两门 carry 得很稳。
}}

\newpage

% ============================================================
\section{Q9: Systematic-Scan Gibbs Sampling}
% ============================================================

\subsection{题目设定}

两个二元变量 $x_1, x_2 \in \{0, 1\}$，联合分布 $\pi$：
\[
  \pi(0,0) = 0.4, \quad \pi(0,1) = 0.1, \quad \pi(1,0) = 0.1, \quad \pi(1,1) = 0.4
\]

直觉：$x_1$ 和 $x_2$ 强烈倾向取相同的值（$P(\text{same}) = 0.8$）。

Systematic-scan Gibbs 的一次 sweep：先从 $\pi(x_1 \mid x_2)$ 更新 $x_1$，再从 $\pi(x_2 \mid x_1)$ 更新 $x_2$。

% ------------------------------------------------------------
\subsection{Part (a): 算 full conditional distributions（4 分）}
% ------------------------------------------------------------

\subsubsection{你的答案}

你算的是 $\pi(x_1 \mid x_2) = \frac{0.1 + 0.4}{0.4 + 0.1 + 0.1 + 0.4} = 0.5$。

问题：这算的是 \textbf{marginal} $\pi(x_1)$，不是 conditional $\pi(x_1 \mid x_2)$。你对 $x_2$ 的所有值求了和（分母是全部四项加起来），但 conditional 应该\textbf{固定} $x_2$ 的值。

\subsubsection{正确做法}

Conditional 的定义：
\[
  \pi(x_1 \mid x_2) = \frac{\pi(x_1, x_2)}{\pi(x_2)}
  = \frac{\pi(x_1, x_2)}{\sum_{x_1'} \pi(x_1', x_2)}
\]
分母只对 $x_1$ 求和，$x_2$ 是\textbf{固定的}。

\textbf{当 $x_2 = 0$ 时}（看 $\pi$ 表里 $x_2 = 0$ 那一列）：
\begin{align}
  \pi(x_2 = 0) &= \pi(0,0) + \pi(1,0) = 0.4 + 0.1 = 0.5 \notag \\[4pt]
  \pi(x_1 = 0 \mid x_2 = 0) &= \frac{\pi(0,0)}{\pi(x_2=0)} = \frac{0.4}{0.5} = \frac{4}{5} \notag \\[4pt]
  \pi(x_1 = 1 \mid x_2 = 0) &= \frac{\pi(1,0)}{\pi(x_2=0)} = \frac{0.1}{0.5} = \frac{1}{5} \notag
\end{align}

\textbf{当 $x_2 = 1$ 时}（看 $x_2 = 1$ 那一列）：
\begin{align}
  \pi(x_2 = 1) &= \pi(0,1) + \pi(1,1) = 0.1 + 0.4 = 0.5 \notag \\[4pt]
  \pi(x_1 = 0 \mid x_2 = 1) &= \frac{0.1}{0.5} = \frac{1}{5} \notag \\[4pt]
  \pi(x_1 = 1 \mid x_2 = 1) &= \frac{0.4}{0.5} = \frac{4}{5} \notag
\end{align}

同理 $\pi(x_2 \mid x_1)$（对称的，因为 $\pi$ 关于 $x_1, x_2$ 对称）：

\[
\pi(x_1 \mid x_2) = \begin{pmatrix} 4/5 & 1/5 \\ 1/5 & 4/5 \end{pmatrix}, \qquad
\pi(x_2 \mid x_1) = \begin{pmatrix} 4/5 & 1/5 \\ 1/5 & 4/5 \end{pmatrix}
\]

其中行是 $x_1$（或 $x_2$）的值，列是 conditioning variable 的值。

\vspace{0.3em}
\noindent\fbox{\parbox{\dimexpr\textwidth-2\fboxsep-2\fboxrule}{%
\textbf{错误根源。} Marginal 是对\textbf{其他变量求和掉}（$\pi(x_1) = \sum_{x_2} \pi(x_1, x_2)$），结果不依赖 $x_2$。Conditional 是\textbf{固定}另一个变量的值再归一化（$\pi(x_1 \mid x_2 = 0) = \pi(x_1, 0) / \pi(x_2 = 0)$），结果随 $x_2$ 的值变化。你在考试中做的是前者，题目问的是后者。
}}

\subsubsection{和 CAVI 的联系}

这个 conditional 就是 Gibbs sampling 每一步用的东西。给定 $x_2$ 的当前值，从 $\pi(x_1 \mid x_2)$ 里\textbf{采样}一个新的 $x_1$。

回忆我们在 L13 学的：CAVI 也是一次更新一个变量，但它是\textbf{确定性的}（设成 conditional 的最优值）。Gibbs sampling 是\textbf{随机的}（从 conditional 里采样）。同一个 conditional，不同用法。

% ------------------------------------------------------------
\subsection{Part (b): 算 $4 \times 4$ transition matrix $P$（4 分）}
% ------------------------------------------------------------

\subsubsection{什么是 transition matrix}

State space 是 $(x_1, x_2) \in \{(0,0), (0,1), (1,0), (1,1)\}$，共 4 个 state。

$P(y \mid x)$ = 从 state $x$ 出发，经过一次完整 sweep（先更新 $x_1$，再更新 $x_2$）后，到达 state $y$ 的概率。

一次 sweep 分两步：
\begin{enumerate}
  \item 固定 $x_2$，从 $\pi(x_1 \mid x_2)$ 采样新的 $x_1$ $\to$ 得到中间 state $(x_1', x_2)$
  \item 固定 $x_1'$，从 $\pi(x_2 \mid x_1')$ 采样新的 $x_2$ $\to$ 得到最终 state $(x_1', x_2')$
\end{enumerate}

所以：$P\big((x_1', x_2') \mid (x_1, x_2)\big) = \sum_{x_1''} \pi(x_1'' \mid x_2) \cdot \pi(x_2' \mid x_1'')$

等等 --- 这里的 $x_1''$ 是中间值。但最终 state 的 $x_1'$ 是 step 1 采出来的，$x_2'$ 是 step 2 采出来的。所以：

\[
  P\big((a, b) \mid (x_1, x_2)\big) = \pi(x_1' = a \mid x_2) \cdot \pi(x_2' = b \mid x_1' = a)
\]

注意：step 2 conditioning on 的是 step 1 \textbf{更新后的} $x_1' = a$，不是原来的 $x_1$。

\subsubsection{逐项计算}

\textbf{从 $(0,0)$ 出发}：$x_2 = 0$

\begin{itemize}[nosep]
  \item $\to (0,0)$: $\pi(0 \mid x_2{=}0) \cdot \pi(0 \mid x_1'{=}0) = \frac{4}{5} \cdot \frac{4}{5} = \frac{16}{25}$
  \item $\to (0,1)$: $\pi(0 \mid 0) \cdot \pi(1 \mid 0) = \frac{4}{5} \cdot \frac{1}{5} = \frac{4}{25}$
  \item $\to (1,0)$: $\pi(1 \mid 0) \cdot \pi(0 \mid 1) = \frac{1}{5} \cdot \frac{1}{5} = \frac{1}{25}$
  \item $\to (1,1)$: $\pi(1 \mid 0) \cdot \pi(1 \mid 1) = \frac{1}{5} \cdot \frac{4}{5} = \frac{4}{25}$
\end{itemize}

Check: $\frac{16+4+1+4}{25} = 1$ \checkmark

由对称性（$\pi$ 在 $0 \leftrightarrow 1$ 下对称），从 $(1,1)$ 出发的行跟从 $(0,0)$ 一样（$0$ 和 $1$ 互换）。类似地可以算从 $(0,1)$ 和 $(1,0)$ 出发的行。

完整的 $P$（行 = 出发 state，列 = 到达 state）：
\[
  P = \frac{1}{25}\begin{pmatrix}
    16 & 4 & 1 & 4 \\
    4 & 1 & 4 & 16 \\
    4 & 16 & 1 & 4 \\
    1 & 4 & 4 & 16
  \end{pmatrix}
\]

\subsubsection{你的答案}

你填了很多 0.25，这对应 uniform transition --- 但实际上因为 $x_1, x_2$ 倾向相同，transition 应该偏向 $(0,0)$ 和 $(1,1)$。

% ------------------------------------------------------------
\subsection{Part (c): 证明 $\pi$ 是 $P$ 的 stationary distribution（4 分）}
% ------------------------------------------------------------

需要验证 $\pi^T P = \pi^T$（即 $\sum_x \pi(x) P(y \mid x) = \pi(y)$ 对所有 $y$）。

你的思路是对的（写了 $\pi^T \cdot T = \pi^T$），但因为 (b) 的 $P$ 算错了，验证也跟着错。用正确的 $P$ 验证：

\[
  \pi^T P = \begin{pmatrix} 0.4 & 0.1 & 0.1 & 0.4 \end{pmatrix}
  \cdot \frac{1}{25}\begin{pmatrix} 16 & 4 & 1 & 4 \\ 4 & 1 & 4 & 16 \\ 4 & 16 & 1 & 4 \\ 1 & 4 & 4 & 16 \end{pmatrix}
  = \begin{pmatrix} 0.4 & 0.1 & 0.1 & 0.4 \end{pmatrix} = \pi^T \;\checkmark
\]

% ------------------------------------------------------------
\subsection{Part (d): Detailed balance 不成立（4 分）}
% ------------------------------------------------------------

Detailed balance 要求 $\pi(x) P(y \mid x) = \pi(y) P(x \mid y)$ 对\textbf{所有} $x, y$ 成立。

只要找一对 $x, y$ 使得等式不成立即可。你在考试中选了 $(0,0)$ 和 $(0,1)$，思路是对的。

\[
  \pi(0,0) \cdot P\big((0,1) \mid (0,0)\big) = 0.4 \cdot \frac{4}{25} = \frac{1.6}{25}
\]
\[
  \pi(0,1) \cdot P\big((0,0) \mid (0,1)\big) = 0.1 \cdot \frac{4}{25} = \frac{0.4}{25}
\]
$\frac{1.6}{25} \neq \frac{0.4}{25}$，所以 detailed balance 不成立。

这不矛盾因为：stationarity（$\pi P = \pi$）只要求\textbf{总流量平衡}（每个 state 的流入 = 流出），不要求每对 state 之间的流量都平衡。Detailed balance 是更强的条件。Systematic-scan Gibbs 满足前者但不满足后者，因为更新顺序是固定的（先 $x_1$ 后 $x_2$），不是随机的。

% ------------------------------------------------------------
\subsection{Part (e): 高维下 Gibbs vs MH（来自 Q10）}
% ------------------------------------------------------------

（留待后续补充。）

% (成绩分析已移至文档开头)

\end{CJK}
\end{document}
